\section*{अध्याय \ref{ch:b1:digestion-absorption} --- पाचन और अवशोषण}

\begin{solution}{exo:b1:digestion-absorption:1}
मुँह: चबाना, लार-ऐमिलेस। ग्रसिका: परिवहन। आमाशय:
अम्ल, पेप्सिन, भंडारण तथा मिश्रण। छोटी आँत: पित्त और अग्न्याशयी
\omterm{def:b1:enzymes:enzyme}{एंजाइम} \omterm{def:b1:digestion-absorption:nutrition}{पाचन} पूरा करते हैं; अवशोषण। बड़ी आँत: जल और
लवण वापस, सूक्ष्मजीवों से \omterm{def:b1:respiration-fermentation:fermentation}{किण्वन}। मलाशय: भंडारण तथा
निष्कासन।
\end{solution}

\begin{solution}{exo:b1:digestion-absorption:2}
आमाशय में pH 2 पर पेप्सिन; छोटी आँत में pH 8 पर अग्न्याशय से आई
ट्रिप्सिन, काइमोट्रिप्सिन तथा कार्बॉक्सीपेप्टिडेस;
और pH 7 पर ब्रश-सीमा की पेप्टिडेस। ज़ाइमोजन के रूप में, ताकि वे उन्हीं
\omterm{def:b1:cell-unit-of-life:cell}{कोशिकाओं} को न पचा डालें जो उन्हें बनाती हैं; वे केवल गुहा में सक्रिय होती हैं।
\end{solution}

\begin{solution}{exo:b1:digestion-absorption:3}
ग्लूकोज: सोडियम सहवहन से भीतर, वाहक से बाहर, प्रवेशिका शिरा और
यकृत तक। \omterm{def:b1:proteins:aminoacid}{ऐमीनो अम्ल}: सोडियम सहवहन से भीतर, वाहक से बाहर, प्रवेशिका शिरा।
\omterm{def:b1:lipids:fattyacid}{वसा अम्ल}: किसी मिसेल से विसरण द्वारा भीतर, फिर एस्टरित, किसी
काइलोमाइक्रॉन में बँधा, और बहिःकोशिकता से दुग्धवाहिनी तथा \omterm{def:b1:mammal-organization:compartments}{लसीका} में।
\end{solution}

\begin{solution}{exo:b1:digestion-absorption:4}
पित्त लवण \omterm{def:b1:lipids:triglyceride}{वसा} को माइक्रोमीटर की बूँदों में पायसित करते हैं और लाइपेस के
उत्पादों को मिसेलों में \omterm{def:b1:digestion-absorption:surface}{आंत्रकोशिकाओं} तक ले जाते हैं। उनके बिना
\omterm{def:b1:lipids:triglyceride}{वसा} बड़ी बूँदों के रूप में रह जाती है जिनका लाइपेस के लिए लगभग कोई पृष्ठ
नहीं, और छूटे \omterm{def:b1:lipids:fattyacid}{वसा अम्ल} झिल्ली तक पहुँचाए नहीं जा सकते: \omterm{def:b1:lipids:triglyceride}{वसा} बिना पची
गुज़र जाती है।
\end{solution}

\begin{solution}{exo:b1:digestion-absorption:5}
ग्लूकोज इकाइयों का $100/162 = \qty{0.62}{mol}$, इसलिए लगभग
\qty{0.62}{mol} बंध और \qty{0.62}{mol} जल (\qty{11}{g});
सोखा गया ग्लूकोज $620/180 = \qty{3.4}{mmol/min}$।
\end{solution}

\begin{solution}{exo:b1:digestion-absorption:6}
अम्ल सीक्रेटिन छोड़ता है, जो अग्न्याशय से ऐसा बाइकार्बोनेट स्रावित
करवाता है जो काइम को उदासीन कर देता है; \omterm{def:b1:lipids:triglyceride}{वसा} और पेप्टाइड कोलीसिस्टोकाइनिन
छोड़ते हैं, जो पित्ताशय को सिकोड़ता है (पित्त घुसता है), अग्न्याशयी
\omterm{def:b1:enzymes:enzyme}{एंजाइमों} को उकसाता है, और आमाशय का खाली होना संदमित करता है, ताकि
ग्रहणी को काइम उतनी ही तेज़ी से मिले जितनी तेज़ी से वह उसे उदासीन तथा पचा
सके।
\end{solution}

\begin{solution}{exo:b1:digestion-absorption:7}
ग्लूकोज अपनी प्रवणता के विरुद्ध आठ गुना सांद्रित हुआ: यानी सक्रिय
परिवहन। सोडियम के बिना परिवहन साम्य पर रुक जाता है: यह
ग्रहण सोडियम-प्रवणता से युग्मित है (सोडियम--ग्लूकोज
\omterm{prop:b1:membranes-transport:secondary}{सहवाहक}), सीधे \omterm{def:b1:water-small-molecules:atp}{ATP} से नहीं।
\end{solution}

\begin{solution}{exo:b1:digestion-absorption:8}
एक माइक्रोमीटर चौड़े काइलोमाइक्रॉन कसी केशिका-दीवार पार नहीं कर सकते
पर खुले सिरों वाली दुग्धवाहिनियों में घुस सकते हैं; \omterm{def:b1:mammal-organization:compartments}{लसीका} रक्त से
गर्दन पर मिलती है, इसलिए \omterm{def:b1:lipids:triglyceride}{वसा} यकृत से पहले समूचे शरीर तक पहुँच जाती है,
जो अन्यथा उसे सबसे पहले पूरा ले लेता। सीधे प्रवेशिका के रक्त में घुसने पर
काइलोमाइक्रॉन यकृत की कोटरिकाएँ बंद कर देते और हर भोजन के बाद \omterm{def:b1:lipids:triglyceride}{वसा}
संसाधित करने की उसकी क्षमता संतृप्त कर देते।
\end{solution}

\begin{solution}{exo:b1:digestion-absorption:9}
वह रोमंथन के लिए रेशे को भिगोती और तैराती है; उसका बाइकार्बोनेट और
फ़ॉस्फ़ेट उन अम्लों का प्रतिरोध करते हैं जो सूक्ष्मजीव बनाते हैं, और रूमेन
को pH 6.5 के पास थामे रखते हैं; और वह रक्त से यूरिया ढोती है, जिसे
सूक्ष्मजीव अमोनिया में और फिर अपने ही \omterm{def:b1:proteins:peptide}{प्रोटीन} में बदल देते हैं, और वह
नाइट्रोजन पुनर्चक्रित कर देते हैं जो अन्यथा मूत्र में खो जाती।
\end{solution}

\begin{solution}{exo:b1:digestion-absorption:10}
अग्न्याशयी ऐमिलेस, प्रोटीन-अपघटनियों, लाइपेस तथा बाइकार्बोनेट के बिना: \omterm{def:b1:carbohydrates:polysaccharide}{मंड}
केवल लार और ब्रश-सीमा से अंशतः पचता है; \omterm{def:b1:proteins:peptide}{प्रोटीन}
पेप्सिन से कटते हैं पर सोखने लायक आकार तक नहीं; \omterm{def:b1:lipids:triglyceride}{वसा} बिलकुल नहीं पचती
और भारी, फीके, चिकने मल में प्रकट होती है। शर्कराएँ और \omterm{def:b1:carbohydrates:glycosidic}{द्विशर्कराएँ}
(ब्रश-सीमा के \omterm{def:b1:enzymes:enzyme}{एंजाइम} बचे रहते हैं), कुछ \omterm{def:b1:carbohydrates:polysaccharide}{मंड}, और छोटे पेप्टाइड फिर भी
सँभाले जा सकते हैं; \omterm{def:b1:lipids:triglyceride}{वसा} और अधिकांश \omterm{def:b1:proteins:peptide}{प्रोटीन} नहीं।
\end{solution}

\begin{solution}{exo:b1:digestion-absorption:11}
अंधनाली विष्ठा वह सूक्ष्मजीवी \omterm{def:b1:proteins:peptide}{प्रोटीन} और वे B विटामिन ढोती है जो
अंधनाल में बनते हैं, जो छोटी आँत से आगे पड़ता है और इसलिए पहली बार गुज़रने
पर पचाया नहीं जा सकता; खाई जाने पर वह आमाशय
और छोटी आँत से गुज़रती है और सोख ली जाती है। गाय आमाशय से पहले \omterm{def:b1:respiration-fermentation:fermentation}{किण्वन}
करती है, इसलिए उसके सूक्ष्मजीव बिना किसी दूसरे गुज़र के आगे पचा लिए जाते
हैं। घोड़े का अंधनाल भी आगे ही है, और घोड़ा न अपनी विष्ठा खाता है न उसके
पास कोई दूसरा रास्ता है: उसका सूक्ष्मजीवी \omterm{def:b1:proteins:peptide}{प्रोटीन} खो जाता है।
\end{solution}

\begin{solution}{exo:b1:digestion-absorption:12}
सूक्ष्मजीव गाय को \omterm{def:b1:carbohydrates:polysaccharide}{सेलुलोस} की ऊर्जा \omterm{def:b1:lipids:fattyacid}{वसा अम्लों} के रूप में देते हैं, गरीब
घास से और यूरिया तक से बना \omterm{def:b1:proteins:peptide}{प्रोटीन}, और हर B विटामिन;
और गाय उन्हें एक गरम, प्रतिरोधित, अवायुजीवी हौज़ देती है जो भरा तथा
हिलता रहता है, और दिन में सौ लीटर लार। इस रचना की लागतें: ऊर्जा का
दसवाँ भाग मेथेन बनकर निकल जाता है, कोई ग्लूकोज नहीं सोखा जाता इसलिए
यकृत को वह सारा प्रोपिओनेट से बनाना पड़ता है, \omterm{def:b1:digestion-absorption:nutrition}{पाचन} में तीन दिन लगते हैं,
और हौज़ तथा उसकी सामग्री जंतु का पाँचवाँ भाग तौलते हैं --- कोई गाय
दौड़ नहीं सकती, और कोई सिंह दौड़ सकता है।
\end{solution}

\begin{solution}{pb:b1:digestion-absorption:1}
\textbf{1.} रेशा \qty{5}{kg}, \omterm{def:b1:carbohydrates:polysaccharide}{मंड} और शर्कराएँ \qty{2}{kg}, \omterm{def:b1:proteins:peptide}{प्रोटीन}
\qty{1.5}{kg}।
\textbf{2.} $0.7\times 5 + 2 = \qty{5.5}{kg}$।
\textbf{3.} $5500\times 17.5 = \qty{96}{MJ}$; VFA $0.75\times 96 =
\qty{72}{MJ}$।
\textbf{4.} $0.08\times 96 = \qty{7.7}{MJ}$; $7700/55 = \qty{140}{g}$
मेथेन, यानी लगभग \qty{200}{L}।
\textbf{5.} $0.15\times 5500 = \qty{825}{g}$।
\textbf{6.} अधिकांश आहारीय \omterm{def:b1:proteins:peptide}{प्रोटीन} रूमेन में सूक्ष्मजीव तोड़ देते हैं और
सूक्ष्मजीवी \omterm{def:b1:proteins:peptide}{प्रोटीन} के रूप में फिर गढ़ देते हैं; सूक्ष्मजीव आगे
एबोमेजम और छोटी आँत तक बहते हैं, जहाँ वे किसी भी मांस की तरह पचा लिए जाते हैं।
\textbf{7.} रेशे का \qty{30}{\%}, यानी \qty{1.5}{kg}: अंशतः
बृहदांत्र में किण्वित, अधिकतर उत्सर्जित।
\textbf{8.} बचा आहारीय \omterm{def:b1:proteins:peptide}{प्रोटीन} $0.4\times 1500 = \qty{600}{g}$,
जिसमें से पचा \qty{360}{g}; सूक्ष्मजीवी \qty{825}{g}; कुल \qty{1185}{g}
$\times 17.5 = \qty{20.7}{MJ}$।
\textbf{9.} $72 + 20.7 = \qty{93}{MJ}$।
\textbf{10.} $72/93 = \qty{77}{\%}$ (पूरे बजटों में \omterm{def:b1:respiration-fermentation:fermentation}{किण्वन} की ऊष्मा और
बृहदांत्र के योगदान की ऊर्जा गिनने पर लगभग \qty{60}{\%}; इस हिसाब से,
तीन-चौथाई)।
\textbf{11.} \qty{93}{MJ} \qty{110}{} से कम है: गाय को अधिक खाना
(\qty{12}{kg}), या बेहतर घास चाहिए, या अपनी \omterm{def:b1:lipids:triglyceride}{वसा} पर हाथ डालना पड़ेगा।
\textbf{12.} प्रोपिओनेट; यकृत में नवग्लूकोजनन।
\textbf{13.} $0.9\times 2000\times 17.5 = \qty{31.5}{MJ}$।
\textbf{14.} $0.7\times 1500\times 17.5 = \qty{18.4}{MJ}$।
\textbf{15.} $0.45\times 5 = \qty{2.25}{kg}$; VFA $0.75\times 2250\times
17.5 = \qty{29.5}{MJ}$।
\textbf{16.} $0.15\times 2250 = \qty{340}{g}$, जो मल में खो जाता है
(पश्चांत्र छोटी आँत से आगे है)।
\textbf{17.} $31.5 + 18.4 + 29.5 = \qty{79}{MJ}$; VFA का हिस्सा $29.5/79 =
\qty{37}{\%}$।
\textbf{18.} नहीं: उसे $110/79\times 10 = \qty{14}{kg}$ घास चाहिए, और
वह दिन में सोलह घंटे खाकर तथा भोजन को दुगनी तेज़ी से गुज़ारकर निभाता है।
\textbf{19.} गाय $93/10 = \qty{9.3}{MJ/kg}$; घोड़ा $79/10 =
\qty{7.9}{MJ/kg}$।
\textbf{20.} सकल ग्रहण $10\,000\times 0.85\times 17.5 =
\qty{149}{MJ}$ (लिग्निन तथा भस्म छोड़कर); \qty{7.7}{MJ} मेथेन उसका
\qty{5}{\%} है; सौ गायें: साल में $100\times 140\times 365 =
\qty{5.1}{t}$ मेथेन।
\textbf{21.} गाय की आँत में $10\times 70/24 = \qty{29}{kg}$ शुष्क पदार्थ
(जोड़ उसका दस गुना जल); घोड़े की $10\times 36/24 =
\qty{15}{kg}$।
\textbf{22.} सिद्धांततः पूरा \qty{340}{g}, जो \qty{6}{MJ} और उसके
आवश्यक \omterm{def:b1:proteins:aminoacid}{ऐमीनो अम्लों} जितना है; किसी घोड़े की अंधनाली सामग्री किसी नरम
पोषक-धनी अंश में उस तरह नहीं बँटती जैसे किसी खरगोश की, और आधे टन का कोई
जंतु ज़मीन से अपनी विष्ठा किसी काम की हालत में वापस नहीं ले सकता।
\textbf{23.} रूमेन के सूक्ष्मजीव यूरिया की नाइट्रोजन और भूसे के कार्बन
से \omterm{def:b1:proteins:peptide}{प्रोटीन} बनाते हैं, और गाय सूक्ष्मजीवों को पचा लेती है: उसे लगभग कोई
आहारीय \omterm{def:b1:proteins:peptide}{प्रोटीन} नहीं चाहिए। घोड़े के सूक्ष्मजीव उसकी छोटी आँत से आगे हैं,
इसलिए उनका \omterm{def:b1:proteins:peptide}{प्रोटीन} खो जाता है, और घोड़े को अपने \omterm{def:b1:proteins:aminoacid}{ऐमीनो अम्ल} भोजन में ही
खोजने पड़ते हैं।
\textbf{24.} घोड़ा ग्लूकोज सोखता है और अवायुजीवी दौड़ के लिए अपनी
पेशियों में \omterm{def:b1:carbohydrates:polysaccharide}{ग्लाइकोजन} संचित करता है; गाय कुछ नहीं सोखती, प्रोपिओनेट से
धीरे-धीरे ग्लूकोज बनाती है, और सौ लीटर का हौज़ ढोती है --- चरने के लिए
गढ़ी हुई, दौड़ने के लिए नहीं।
\textbf{25.} इस बजट पर गाय की उपापचेय ऊर्जा का लगभग तीन-चौथाई
(पूरे हिसाब में \qtyrange{60}{70}{\%}) वाष्पशील
\omterm{def:b1:lipids:fattyacid}{वसा अम्लों} के रूप में, जबकि घोड़े का एक-तिहाई; और प्रति किलोग्राम घास
\qty{7.9}{MJ} के सामने \qty{9.3}{MJ}।
\end{solution}
